\subsection{Confirmation experiment with three repetitions}
After the pilot, the same 39 graphs, three $(h,k)$ pairs, two backends, and
30-second budget per run were retained, with three repetitions per backend
and instance. The file \path{results/runs/general_confirm_r1.csv} contains
\ConfirmObservations{} observations for \ConfirmInstances{} graph/$(h,k)$
instances. These data are analyzed separately from the pilot; additional
repetitions do not increase the number of independent graphs.
Every witness was checked using shortest-path distances on the regenerated
graph. No inconsistency in spans, bounds, or OPT values was found across
repetitions or backends. Exact reference formulas cover
\ConfirmReferenceInstances{} instances, with \ConfirmReferenceMatches{}
matching observations. This is neither independent UNSAT-certificate
verification nor a reproduction of historical timing measurements.

CaDiCaL records \ConfirmSATOpt/351 OPT observations within budget and Gurobi
\ConfirmILPOpt/351, leaving \ConfirmFeasible{} FEASIBLE observations across
both backends. There are \ConfirmOutsideBudget{} observations reporting OPT
outside the external budget. Requiring OPT in all three repetitions yields
\ConfirmSATStable/117 instances for CaDiCaL and \ConfirmILPStable/117 for
Gurobi, with \ConfirmJointStable{} instances satisfying the criterion for
both. Repeated-run variability is retained rather than reporting the
fastest run as the representative time.

\begin{table}[htbp]
\centering\small
\caption{Confirmation results: distinct instances, OPT observations within
30 seconds, and instances OPT in all three repetitions. ER groups are
separated by edge probability $p$.}
\label{tab:confirmation-coverage}
\input{\ConfirmReportDir/summary.tex}
\end{table}

\begin{figure}[htbp]
\centering
\includegraphics[width=\linewidth]{\ConfirmReportDir/coverage.pdf}
\caption{Instances solved to optimality in all three repetitions, by group.
The denominator is the number of graph/$(h,k)$ instances in the group,
not the number of repetitions.}
\label{fig:confirmation-coverage}
\end{figure}

\begin{table}[htbp]
\centering\small
\caption{Runtime on the common subset solved to optimality in all three
repetitions by both backends. The median of three runs is computed per
instance, followed by the median across instances in each group. A dash
indicates no eligible instances, not zero seconds.}
\label{tab:confirmation-time}
\input{\ConfirmReportDir/timing.tex}
\end{table}

\begin{figure}[htbp]
\centering
\includegraphics[width=\linewidth]{\ConfirmReportDir/runtime.pdf}
\caption{Each point represents an instance in the common all-three-OPT
subset. Coordinates are median runtimes; error bars show the minimum and
maximum of three runs, not confidence intervals. Points below the diagonal
indicate lower SAT runtime. Interpret together with coverage in
Figure~\ref{fig:confirmation-coverage}.}
\label{fig:confirmation-runtime}
\end{figure}

These runtime medians are subject to selection bias because they include
only the common solved subset; they do not represent instances remaining
FEASIBLE. OPT-observation counts and the all-three-OPT criterion provide
complementary information about optimality coverage within budget. The
results establish neither general superiority of one formulation nor any
comparison with GA. Tables and figures are exported with data and audit-source
hashes. The preceding model-size counts remain a separate common-span
analysis, not counts of models used in the confirmation runs.
\FloatBarrier
